\chapter{Know What the Number Means}
\chaptermark{Know What the Number Means}

A large financial number is an unfinished sentence.

In Miami, an industrial-property developer described leaving a long apprenticeship
to start a company with his partner. They mortgaged their homes, drew on credit
lines and cards, and put everything into the business. It succeeded. In Austin,
Matt Teifke described being down several million dollars on one shareholding
while remaining a large owner. In New York, a trader remembered margin calls he
managed to cover and the sleep they cost him.

Each account contains a dramatic quantity, even where no amount was stated:
everything exposed, millions of marked loss, cash demanded now. They describe
different economic states. The first is a set of claims placed behind a company
before its outcome was known. The second is a change in the quoted value of an
owned position. The third is another party's right to require cash. Blending
them into one lesson about courage would teach almost nothing.

The judgment about how much risk a person can bear comes later. First we need
the grammar. What does the number measure? To whom does it belong? When is it
true? What can its holder actually do with it, and who has a prior claim?

\section{Complete the address before reading the amount}

An Austin lawyer gives us the cleanest demonstration. Asked for the most money
he had made in a year, he answered with what he called paper money: a reported
\$17.5 million value attached to stock options in companies going public near
the height of the dot-com boom. The shares were subject to a six-month lockup.
By the time he says he could act, the position was worth about \$750,000; he
recalls eventual sale proceeds of roughly \$350,000.

The three amounts concern the same stake at three moments, but they do not answer
the same question. The first was neither annual income nor cash available for
use. The second was a later market value. The third was reported sale proceeds,
still before whatever tax, fees, and personal obligations would determine usable
cash. His experience led him to favor selling early. That is an understandable
conclusion from his path, not a rule that settles every concentrated holding.
What survives the anecdote is the need to state the conditions attached to a
price.

A complete number address has eight lines:

\enlargethispage{\baselineskip}

\begin{description}[leftmargin=0em,style=unboxed,itemsep=0.1em,
  topsep=0.2em,parsep=0em,labelsep=0.55em]
\item[Entity and share] Name the person, household, company, fund, asset, or
  transaction, then the exact ownership interest to which the number applies.
\item[Time] Give the period for a flow or the date for a balance and mark. Say
  whether it is historical, annualized, forecast, cumulative, or as of now.
\item[Category] Distinguish revenue, contribution, operating profit, net income,
  cash flow, salary, distribution, debt, enterprise value, equity value, and
  personal available resources.
\item[Scope] Mark gross or net, consolidated or partial, before or after tax,
  and the fraction attributable to the stated owner.
\item[State] Mark collected or receivable, realized or estimated, liquid or
  illiquid, unrestricted, pledged, locked, contingent, or transferable.
\item[Claims ahead] Record delivery duties, payroll, tax, debt, senior rights,
  reserves, guarantees, and any condition that stands between the amount and use.
\item[Source and confidence] Separate a bank record, account, contract, appraisal,
  or transaction statement from a dashboard, estimate, interview, or memory.
  Use a range when the evidence supports only a range.
\item[Decision] State what choice the number is meant to support and what happens
  if its category, timing, or confidence is wrong.
\end{description}

This is practical rather than ceremonial. A company may have strong revenue and
weak collection. A profitable owner may receive no distribution. A valuable
shareholding may be inaccessible when a household needs cash. The missing noun
often changes the decision more than the amount does.

\section{Follow the customer all the way to cash}

Revenue is the number most often hurried into an owner's life. The operating
record from Part III gives it a route:

\begin{equation}
\begin{gathered}
\text{attention}\rightarrow\text{qualified decision}\rightarrow\text{exchange}\\
\rightarrow\text{delivery}\rightarrow\text{acceptance}\rightarrow\text{collection}.
\end{gathered}
\end{equation}

A signed order may create a duty before it creates revenue. Recognized revenue
may precede collection. Collected cash may arrive with refund, warranty,
commission, support, or recovery obligations. Only after the direct burden of
the exchange is counted do we see its contribution:

\begin{equation}
\ifdim\paperwidth<450pt
\begin{gathered}
\text{collected contribution}\\
=\text{collected cash}-\text{refunds}\\
\phantom{=}-\text{direct burden}.
\end{gathered}
\else
\text{collected cash}-\text{refunds and direct burden}
=\text{collected contribution}.
\fi
\end{equation}

That contribution must join the rest of the company rather than leap into the
owner's wallet. Operating costs, working-capital needs, equipment, debt service,
tax, repair, and protected reserves still have dates. What remains is available
company cash. A distribution becomes personal cash only after authority,
documents, law, tax, and the company's existing promises permit it. One dollar
cannot be counted both as the company's survival reserve and as household
spending.

The Scottsdale interviews supply several useful addresses. A healthcare founder
reported about \$100 million of annual company revenue, sole ownership, and an
estimated personal net-worth range of \$500 million to \$700 million. The host
later described cash as having been retained in the company. These are four
claims: company revenue for a period, an ownership right, a personal estimate at
an unstated date, and company liquidity assigned to operations. None converts
the other three into spendable personal cash.

An online fitness founder reported that the current year would exceed \$20
million of annual recurring revenue. That statement needs both a time label and
a recurrence test. Is the amount contracted, collected, or annualized from a
recent period? Do customers renew voluntarily? What do fulfillment, support,
refunds, acquisition, and recovery cost? Recurrence becomes economically useful
when continuing value earns continuing payment, not when a dashboard repeats a
monthly figure twelve times.

An industrial distributor reported about \$250 million of revenue in strong
years. A later operator reported \$11 million across businesses and recommended
an eighteen-month cash-flow model. The first two are historical company flows.
The last is a forecast horizon. It is not a claim that eighteen months of cash
already sits in an account. A model becomes a runway only after dated liquid
resources and a credible stressed outflow have been supplied.

\section{Reconcile before you admire}

A Houston restaurant owner offered three views of one operation: a
4,600-square-foot restaurant, a one-million-dollar December, and annual revenue
just under \$10 million. He also estimated annual revenue density at roughly
\$2,300 to \$2,400 per square foot.

The statements do not reconcile exactly. Ten million dollars divided by 4,600
square feet is about \$2,174 per square foot. Repeating December for twelve
months would produce \$12 million. The quoted density implies yet another annual
range. Seasonality, rounding, usable versus total area, and different reporting
periods could explain the gaps. Choosing among them without records would turn
an honest unknown into a fabricated answer.

The operation may still achieve unusually high sales from a compact room. Keep
the reported observations, label the mismatch, and recover the numerator,
denominator, period, and records before comparing the ratio. Then ask what the
ratio leaves outside: food and labor, rent, waste, payment fees, insurance,
equipment, tax, owner labor, employee bonuses, and the service intensity needed
to know the customer at every table.

This is the reconciliation habit. When a peak month and a year disagree, do not
annualize whichever looks better. When a ratio is offered, recover both sides.
When a total spans several companies, leave it at the group until an allocation
is supported. When a profit question receives a revenue answer, preserve the
answer and correct the category. A discrepancy is an invitation to inquire, not
permission either to accuse or to invent.

\section{A valuation is a proposition}

In another Houston interview, a private-equity operator was asked about the most
money made in a year. He redirected the question away from dividends and company
earnings to a reported portfolio range of about \$420 million. The mismatch is
the useful part. Income is a flow during a period; portfolio value is a stock
estimated at a date.

For a public security, a visible quote still leaves selling restrictions, tax,
concentration, market depth, and movement between quotation and execution. A
private mark may also depend on an old financing round, a different class of
shares, liquidation preferences, dilution, debt, transfer limits, sparse trades,
and a hoped-for exit. A portfolio total can join positions measured on
incompatible dates and terms. It remains a proposition about realizable value,
not a bank balance.

For an immediate decision by person or entity \(p\), use a narrower quantity:

\begin{equation}
A_{p,t}=C^{\mathrm{free}}_{p,t}+L^{\mathrm{reliable}}_{p,t}
       -O^{\mathrm{due}}_{p,t}-R^{\mathrm{protected}}_{p,t}.
\end{equation}

Here unrestricted cash and reliably realizable liquid assets are reduced by
obligations due and reserves already assigned to tax, payroll, health, repair,
or survival. Illiquid and restricted holdings remain visible on a separate
balance sheet. The expression is a decision aid, not an accounting, legal, tax,
or investment definition.

Now the Austin paper loss can be read without pretending to settle the
investment. Teifke reported being down several million dollars on one cannabis
shareholding while remaining a large shareholder. He cited the company's
quarterly revenue and market value as reasons for continued belief. Those
figures concern the company; his loss concerns his position. To connect them we
would still need the valuation date and category, debt and senior claims, share
class and count, acquisition basis, dilution, liquidity, ownership fraction,
and evidence that the reported revenue could become durable owner cash. A low
ratio can reflect opportunity, deterioration, missing claims, weak evidence, or
some combination. The interview alone cannot choose.

The industrial-property success also belongs on a different line. Mortgaged
homes, credit lines, and cards describe the financing and household exposure;
the successful company describes the later outcome. Without amounts, rates,
recourse, ownership, cash timing, and the household floor, the account cannot be
used as an underwriting rule. Success does not rewrite the address that existed
before success was known.

The New York margin call changes the state again. The trader reported a winning
Amazon trade, later margin calls, lost sleep, and an ability to provide the cash
required. No amount was given. The call matters because a lender had a dated
right, not because the quoted loss was large or small. Covering it records the
outcome of that episode; it does not tell another household what exposure to
accept. Once the claims are named, the next chapters can judge the debt, loss
capacity, and timing without confusing a mark with a demand.

\section{Do not let a peak year become a permanent burden}

A Houston luxury concierge answered a best-year question with a range of
\$2.5 million to \$5 million. The exchange does not establish whether the range
was personal income, company revenue, profit, or cash after tax. His durable
distinction was between earning money and retaining it. A large flow can pass
through a person's life without creating resilience.

An anonymous commercial-real-estate broker made the distinction through a
purchase he would reconsider. He reported a seven-figure year and warned his
younger self against immediately answering it with a conspicuous sports car.
The car is not financially wrong by definition. The danger lies in allowing an
exceptional flow to establish recurring housing, transport, debt, insurance,
maintenance, and social expectations that ordinary years must carry.

Before accepting a permanent commitment, calculate from personal cash that has
been collected, taxed appropriately, and released from prior duties. Price the
continuing cost as well as the purchase. Then run the household through a weak
year, an interruption, and a delayed recovery. The expense may still belong in
the desired life; it should not be needed to prove that a peak once happened.

Visible restraint proves little by itself. A person may spend modestly from
contentment, fear, debt, inaccessible assets, or obligations no outsider can
see. What matters is whether resources and commitments can continue to support
the chosen life.

\clearpage

\section{Open Record 4 with one addressed number}

Choose the number now exerting the most influence over a business or household:
a revenue target, valuation, account balance, debt, distribution, purchase, or
expected sale. Add one page to the wealth notebook. Write \emph{unknown} where
the evidence ends.

\begin{enumerate}[itemsep=0.12em]
\item Copy the claim exactly and record its source and date.
\item Name the entity, owner, ownership share, period or valuation date, and
  category.
\item Mark gross or net, company or personal, cash or accrual, historical,
  annualized, forecast, realized, or estimated.
\item Record liquidity, restrictions, collateral, transfer rights, and every
  condition that can delay or change use.
\item List delivery duties, working capital, tax, debt, senior claims,
  guarantees, and protected reserves standing ahead of the decision.
\item Reconcile the claim to Record 3: attention, decision, exchange, delivery,
  acceptance, collection, contribution, company cash, financing, and any
  authorized owner distribution.
\item Reconcile related statements: month to year, revenue to contribution,
  profit to cash, enterprise value to owner proceeds, and portfolio mark to
  realizable liquidity.
\item Give the result a range and confidence appropriate to the evidence. Name
  the decision it supports, the consequence of error, and the next fact that
  would materially change it.
\end{enumerate}

Do this before admiration, comparison, borrowing, spending, or advice. Some
fortunes will become smaller; some modest balances will become more dependable;
some businesses will reveal that their strongest number cannot yet carry their
next promise.

Once a claim has an address, debt can be judged by the cash and asset expected
to carry it. The lender may change while the obligation remains.
